\begin{longtable}{@{}>{\raggedright\arraybackslash}p{0.34\textwidth} >{\raggedright\arraybackslash}p{0.58\textwidth}@{}}
\caption{The methods of the studio and the acceptance criterion the design document fixes for each: the published methods and the methods proposed here.}\label{tab:ladder}\\
\toprule Method & Acceptance criterion\\\midrule\endfirsthead
\toprule Method & Acceptance criterion\\\midrule\endhead
\multicolumn{2}{@{}l}{\textit{Published methods}}\\
Superposition of prisms and lenses & \footnotesize on the 10x10 task with a flat beam, efficiency within 0.10 of 0.29 and uniformity below 0.05\\
Random superposition & \footnotesize on the 10x10 task with a flat beam, efficiency above superposition of prisms and lenses and at least 0.65; uniformity below 0.35 (DinHot measured 0.25 at 256 squared against 0.01 at 768 squared in the source; the ordering, not the value, transfers)\\
Random mask & \footnotesize efficiency about 1/M; at the source's 768 squared grid, uniformity within the spread of 0.58\\
Gerchberg-Saxton, point-trap and FFT forms & \footnotesize the two forms agree to $10^{-9}$ rad; 30 iterations on the 10x10 task with a flat beam reach efficiency at least 0.9 and uniformity near 0.6\\
Generalized adaptive additive & \footnotesize uniformity above GS at equal iterations on the 10x10 task\\
Weighted GS & \footnotesize efficiency and uniformity at least 0.9 within a median of 15 iterations on the 10x10 task with a flat beam (the source reports about 10); the 3x3x3 lattice above 0.9 and 0.98\\
Phase-fixed weighted GS & \footnotesize fewer iterations than weighted GS to a relative standard deviation of 0.5 percent on a 50x30 array, with the phase fixed after 60 iterations in the acceptance run (the code default `n\_fix` is 12, as in Kim et al. 2019; at 12 it stalls at 3.58 percent on DinHot's grid, re-measured 2026-10-30); from the same start phase-fixed weighted GS matches slmsuite 0.5.0's WGS-Kim to 1e-4 percent (the first field is not reweighted, g\_0 = 1)\\
MRAF & \footnotesize on disc, ring and bar targets, signal-region error at least five times below GS at the cost of efficiency\\
Direct search & \footnotesize uniformity at least 0.98 on the 10x10 task; under 5 percent rms variation on a quasi-periodic set\\
Conjugate-gradient minimization & \footnotesize the analytic gradient matches finite differences; a square lattice below 1 percent error\\
Wirtinger gradient descent with Adam & \footnotesize lower trap-intensity error than GS at equal iterations on the 10x10 task\\
Closed-loop WGS on an emulated camera & \footnotesize from a 19 percent rms open-loop spread, at most 2.1 percent in 20 iterations at gain 0.7; the acceptance runs on the emulated camera (`core/camera.py`), which reproduces the published feedback curves (Nogrette 2014: 19 to 1.4 percent in 20 iterations; Kim 2019: 22 percent to 1.59 plus or minus 0.18 percent in three phase-fixed steps over 8 ensembles, 1.4 in the histogram of one run; the bounds 2.1 and 1.8 are set above them)\\
Device-aware CG with pixel crosstalk & \footnotesize on the crosstalk device, lower error than conjugate-gradient minimization without the model on a large disc; a staircase line search lets it run its budget on a quantizing device, where its scored spread does not improve systematically (a measured negative)\\
Zernike correction and sensorless modal recovery & \footnotesize recovers an injected set of Noll modes 4 to 11 within 5 percent of their amplitude\\
Laguerre-Gaussian vortex traps & \footnotesize the ring phase winds by 2 pi l with an error under 0.1 rad\\
Hungarian assignment and linear phase interpolation & \footnotesize no crossing moves for alpha 1.5 on a 7x7 to 3x3 sort; a moving trap above 70 percent of its mean intensity on the frames themselves (the minimum over the frames of the sequence over their mean; between frames, under the blend model of flicker-bounded phase-continuous sequence's transient, it is not: notebook 07-04)\\
Multi-plane CNN hologram with an optional volumetric GS refinement (DeepCGH-style) & \footnotesize on 162 held-out three-plane volumes, the hybrid (the network's phase, refined by volumetric 3D GS inside the same wall-clock budget only when that raises the volume accuracy) reaches a higher accuracy than volumetric 3D GS alone at that budget (measured median ratio 1.140); the refinement was kept on none of the volumes, so the measured method is the network alone\\
Optimal-transport initialized MRAF & \footnotesize lower error than MRAF from a random start on the shaped targets\\
Multi-plane GS & \footnotesize on a three-plane target, every plane mean within 25 percent of the overall mean (volumetric 3D GS of Whyte and Courtial)\\
\multicolumn{2}{@{}l}{\textit{Methods proposed here}}\\
Device-aware conjugate-gradient refinement & \footnotesize on the crosstalk device, relative standard deviation at least two times below phase-fixed weighted GS, efficiency at most 10 percent lower; measured in the 0.02.004 device-aware paper (10x10 task, simulated, five seeds, the weight update off as the default now is): a geometric mean of 201.3 times (95 percent interval 132.9 to 304.9) at 0.999 of the efficiency on the crosstalk device, 23.5 times (18.6 to 29.6) at 0.997 on the 4-bit device with crosstalk, fill factor and flicker, and 2.0 times on the full C08 device (the bake, three seeds: 1.906, 1.997 and 2.228); against the closest published method the result is negative: the normalized-error loss of Schroff et al. 2023 through the same warm start, device model and conjugate gradient gives 159.9 times (107 to 239.1) at 0.999 on the crosstalk device and 26.1 times (20.2 to 33.9) at 0.997 on the 4-bit device, with overlapping intervals, so the method's own elements (the efficiency-level loss, the weight update, the best-iterate safeguard) add nothing measurable over it, and the weight update, off by default since 0.02.004, raises the spread 7 times on the crosstalk device (24.4 times, 17.5 to 34.1, against 201.3); the gain belongs to the device model in the loop (without it the factor is 1.0); in the benchmark, on the C08 device at 100 iterations with the default parameters, device-aware conjugate-gradient refinement meets both thresholds on 4 of 6 geometries (G2 ring of 36 traps 2.04 times, G3 random 100 traps 4.82, G6a 2.15, G6b 2.08; median factor 2.06 from 1.05 to 4.82 at a median efficiency ratio of 0.969) and not on G1 square 10 x 10 (1.90) or G4 quasi-periodic 119 (1.05); device-aware conjugate-gradient refinement's cells hold 2 seeds; a hologram takes 34 s on the crosstalk device (59 s on the 4-bit device)\\
Flicker-bounded phase-continuous sequence & \footnotesize minimum transient intensity of a moving trap at least 0.9 of its mean, against LPI on the same assignment and against the per-frame GS with phase induction of Kim et al. 2019: on the acceptance task (13 frames, 0.31 Fourier units per frame, the six seeds of the 0.02.004 flicker-bounded paper) flicker-bounded phase-continuous sequence holds 0.915 to 0.932 (median 0.923) against LPI 0.551 to 0.709, so the criterion holds against LPI, and the per-frame GS reaches 0.918 to 0.965 at five iterations per frame and 0.880 to 0.966 at ten, so against it the result is a trade and not an improvement: flicker-bounded phase-continuous sequence's largest phase step between frames is smaller (0.077 rad against 0.408 rad) and its per-frame spread lower, the per-frame method keeps the higher transient minimum, and the gate test asserts at least 0.9 on its seeds 1 to 3 (0.916, 0.920 and 0.926 against LPI 0.709, 0.614 and 0.551); with the frames that keep the peak step per frame at most 0.2 Fourier units (a bound set on C10's sort) long lattice moves reach medians of 0.933 and 0.936, 6 of 8 seeds each, and flicker-bounded phase-continuous sequence is above LPI on 7 of 8 seeds of each (above the per-frame GS at five iterations on 2 and on 1 of 8); on C10's sort of 16 movers (six seeds of the paper) flicker-bounded phase-continuous sequence needs the end phases and 137 frames to reach 0.909 to 0.917 (median 0.914) on 6 of 6 seeds, against LPI 0.044 to 0.738 and the per-frame GS 0.922 to 0.956 at five iterations (flicker-bounded phase-continuous sequence above it on 0 of 6), and at 25 frames no method reaches 0.9: flicker-bounded phase-continuous sequence is at most 0.412, LPI is above flicker-bounded phase-continuous sequence on 4 of 6 seeds and both per-frame sequences are above both of them; the model is DinHot's own and is not calibrated against a measured display\\
Learned warm start for phase-fixed WGS (refuted) & \footnotesize measured negative: on held-out geometries the network start needs a median 8 iterations to 1 percent against 14 from a random start (1.75 times, below the 2 times claimed), and a control with the same early phase fix and no network gains as much; the model is not released and the claim learned warm start is reported as refuted\\
Aberration-robust hologram & \footnotesize over a 200-sample ensemble of residual aberrations, the 95th percentile of the relative standard deviation at least 1.25 times below phase-fixed weighted GS, efficiency at most 10 percent lower (measured, 0.02.004: a median of 1.435 times, 1.42 to 1.57 over three seeds, at 0.982 of the efficiency; over ten seeds a median of 1.439 times (bootstrap interval 1.38 to 1.48, the statistic of the aberration-robust paper), at 0.982, 0.981 to 0.983, with every seed's point estimate above 1.25 and the bootstrap interval of a single 200-sample ensemble above it on 8 of 10 seeds); robustness to pointing errors is measured and reported, not claimed (1.04 times at the acceptance budget; 1.48 times at 0.10 bins and 0.912 of the efficiency with a displacement prior alone and five times the budget)\\
Efficiency-uniformity continuation & \footnotesize hypervolume at least that of the GS, generalized adaptive additive GS, direct search and weighted GS points on the 10x10 task (measured, 0.02.004, eight seeds of the continuation's paper: the headline is the union of the five published points, those four and the phase-fixed WGS, phase-fixed weighted GS: 0.8223 for the continuation against 0.8104, a paired difference of 0.0119 with a 95 percent interval of 0.0097 to 0.0141, positive on 8 of 8 seeds; the union of the four points of the criterion gives 0.8071 and a difference of 0.0152 (0.0125 to 0.0179), 8 of 8); the gain is small and shrinks at a stricter reference point: at the efficiency floors (0.5, 0) and (0.7, 0) the difference is 0.0117 and 0.0113, 97.9 and 95.2 percent of the headline, and at the uniformity floors (0, 0.9) and (0, 0.95) it is 0.0015 (0.0008 to 0.0022) and 0.0011 (0.0005 to 0.0018), 12.4 and 9.6 percent of it (31.3 and 29.1 percent against the four points), because the area the curve adds lies where the uniformity is far below one; in the benchmark (0.02.004 run, 20 seeds) the stored continuation curve (efficiency-uniformity continuation's converged run of the first seed) reaches that hypervolume on 3 of the 7 geometries with a curve against the medians at 100 iterations on the ideal device, losing on G1 (0.836 against 0.839), G4 (0.771 against 0.772), G6a (0.791 against 0.813) and G6b (0.725 against 0.748), by at most 0.023, and on 5 of 7 against those at 30 iterations\\
\bottomrule\end{longtable}
